\begin{helper}[Controlled interchange gives a repeated suffix color]
Let \((\varphi,T)\) be a minimum bad pair with
\(2<\rho(T)<|E(T)|\).  Let \(i<j<|E(T)|\), with
\(\rho(T)\le j\), and let distinct colors \(\alpha,\beta\) be missing
at \(y_i,y_j\), respectively.  If both colors are unused through
\(y_j\) and \(\alpha\) is also missing at the final vertex \(y_n\),
then this configuration is impossible.
\end{helper}
\begin{proof}
Unfold \noderef{TashTreeMinimumBadPair} and take its bad-pair
conjunct.  By \noderef{TashTreeBadPair} this conjunct contains a
\(k\)-triple, and \noderef{KTriple} says in particular that
\(\varphi\) is proper.  We use this properness below for every
incidence-graph endpoint, spanning-path, and swap-missing assertion.
Also, by \noderef{TreeColorUnusedThrough}, the two unused-through
hypotheses say exactly that every displayed edge \(e_r\),
\(2\le r\le j\), has color different from \(\alpha\) and from
\(\beta\).

Apply \noderef{TashTreeControlledInterchange} to \(\alpha,\beta\) at
\(y_i,y_j\), using those edgewise nonuse assertions.  Write \(Q\) for the returned
\((\alpha,\beta)\)-component and \(\varphi'\) for its component swap.
That result says that \(Q\) is a path component containing both
\(y_i\) and \(y_j\), that these two distinct vertices have at most one
neighbor in the two-color incidence graph, that \(\alpha\) is missing
at \(y_j\) under \(\varphi'\), and that \((\varphi',T)\) is again a
minimum bad pair.

We claim that \(y_n\notin Q\).  The inequalities
\(i<j<n\), together with the displayed-vertex injectivity in
\noderef{TreeSequence}, make \(y_i,y_j,y_n\) pairwise distinct.  Since
\(\alpha\) is missing at \(y_n\), the properness extracted above and
\noderef{MissingColorIncidenceEndpoint} says that \(y_n\) also has at
most one neighbor in the two-color incidence graph.  On the other hand,
using the same properness, \noderef{TwoColorComponentSpanningPath}
realizes all of \(Q\) as the
support of one simple path.  A nontrivial simple path has exactly two
vertices of degree at most one in its support: here they are the two
distinct vertices \(y_i\) and \(y_j\).  Thus a third such vertex
\(y_n\) cannot belong to \(Q\), proving the claim.

The outside-component clause of
\noderef{TwoColorComponentSwapMissingBehavior} now preserves the whole
missing-color set at \(y_n\); its properness premise is precisely the
properness derived from \noderef{KTriple}.  Hence \(\alpha\) is missing under
\(\varphi'\) both at \(y_j\) and at \(y_n\).  The unchanged tree has
\(\rho(T)\le j<n\), so these witnesses meet exactly the repeated suffix
hypothesis of \noderef{TashTreeCaseRepeatedOnPath} for the returned
minimum bad pair.  That result gives a contradiction.
\end{proof}
